\documentclass[aspectratio=169]{beamer}

\usepackage[T1]{fontenc}
\usepackage[utf8]{inputenc}
\usepackage{lmodern}
\usepackage{hyperref}

\usetheme{Madrid}
\setbeamertemplate{navigation symbols}{}

\title[Research Progress Update]{Research Progress Update:\\Profiling AI Runtime Categories for Shared-RAN Inference Runtime Selection}
\subtitle[Runtime-system selection]{Category-level profiling and runtime-system selection for edge AI inference on shared RAN infrastructure}
\author{}
\date{}

\newcommand{\slidecite}[1]{\vspace{0.35em}\par{\tiny\raggedright\textit{Citations: }#1\par}}

\begin{document}

\begin{frame}
  \titlepage
\end{frame}

\begin{frame}[t,shrink=10]{Slide 1. What I Have Been Working On}
\footnotesize
\begin{itemize}
  \item I am profiling \textbf{AI inference runtime categories} to understand which ones are technically suitable for shared RAN-edge execution.
  \item The runtime categories I am studying are:
  \begin{itemize}
    \item Transformer / autoregressive (\textbf{Llama 3.2 3B + llama.cpp}),
    \item Diffusion (\textbf{SD1.5 + LCM + TensorRT}),
    \item Computer vision (\textbf{RTMPose-m + TensorRT}),
    \item Voice-related (\textbf{Emformer-style streaming ASR + ONNX Runtime Mobile}),
    \item Translation (\textbf{M2M100-418M + CTranslate2}),
    \item Volume rendering (\textbf{Instant-NGP + tiny-cuda-nn}).
  \end{itemize}
  \item For each category, the profiling focus is not just the model architecture, but the \textbf{runtime engine and serving stack} used to execute the task.
  \item Each category is profiled against a runtime suitability scorecard.
\end{itemize}
\slidecite{Polese et al. (2025); Feng et al. (2026); profile\_doc.md packet summaries.}
\end{frame}

\begin{frame}[t,shrink=6]{Slide 2. Suitability Scorecard}
\footnotesize
\begin{itemize}
  \item \textbf{Preemption resilience}
  \begin{itemize}
    \item Meaning: how safely execution can be interrupted and later resumed.
    \item Impact: determines how much work is lost when resources are reclaimed.
  \end{itemize}
  \item \textbf{Micro-segmentation}
  \begin{itemize}
    \item Meaning: the smallest execution unit the runtime can expose, such as a token step, operator group, tile, ray batch, or denoising step.
    \item Impact: determines whether work can fit short execution windows.
  \end{itemize}
  \item \textbf{State parking}
  \begin{itemize}
    \item Meaning: what runtime state must be saved when compute disappears, and how easily that state can be restored.
    \item Impact: determines pause/resume cost and whether progress survives execution gaps.
  \end{itemize}
  \item \textbf{Tight VRAM compliance}
  \begin{itemize}
    \item Meaning: whether the runtime can stay stable under a strict GPU-memory cap.
    \item Impact: determines OOM risk, admission stability, and multi-tenant feasibility.
  \end{itemize}
  \item \textbf{Throughput}
  \begin{itemize}
    \item Meaning: how much useful work the runtime can sustain under a given GPU share or lease envelope.
    \item Impact: determines whether the runtime remains practical when GPU availability is reduced.
  \end{itemize}
\end{itemize}
\slidecite{Jin et al. (2025); USHER (2024); Orion (2024); RAVAS (2023).}
\end{frame}

\begin{frame}[t,shrink=8]{Slide 3. How Do I Want to Profile?}
\footnotesize
\begin{itemize}
  \item I want to profile the \textbf{runtime system}, not only the model family.
  \item The first profiling strategy is:
  \begin{itemize}
    \item run the runtime under different GPU-availability configurations inspired by \textbf{Weaver-style profiling baselines},
    \item start with \textbf{full GPU availability},
    \item then test \textbf{reduced or shared GPU availability} such as \textbf{MPS-style sharing},
    \item and profile the runtime behavior against the scorecard under each configuration.
  \end{itemize}
  \item For each configuration, I want to observe:
  \begin{itemize}
    \item preemption resilience,
    \item micro-segmentation,
    \item state parking,
    \item tight VRAM compliance,
    \item throughput.
  \end{itemize}
  \item The purpose is to see which runtime-system stacks remain viable as GPU availability changes, instead of only measuring them at full-device conditions.
\end{itemize}
\slidecite{Weaver workload profiling note; USHER (2024); Orion (2024); vLLM/PagedAttention (2023).}
\end{frame}

\begin{frame}[t]{Slide 4. Transformer / Autoregressive}
\small
\begin{itemize}
  \item Category type: streaming, stateful generation with a growing \textbf{KV cache}.
  \item Representative packet-1 anchor: \textbf{Llama 3.2 3B + llama.cpp}.
  \item Current default safe boundary: \textbf{token-step}.
  \item Main profiling question: can the runtime support useful progress, bounded memory growth, and pause/resume around decoder-side state?
  \item Imported mechanism evidence:
  \begin{itemize}
    \item \textbf{vLLM / PagedAttention},
    \item \textbf{Orca},
    \item \textbf{DistServe},
    \item \textbf{SpotServe},
    \item \textbf{CacheGen}.
  \end{itemize}
\end{itemize}
\slidecite{PagedAttention / vLLM (2023); Orca (2022); DistServe (2024); SpotServe (2024); CacheGen (2024).}
\end{frame}

\begin{frame}[t]{Slide 5. Diffusion}
\small
\begin{itemize}
  \item Category type: one-shot \textbf{latent diffusion} with iterative denoising steps.
  \item Representative packet-1 anchor: \textbf{SD1.5 + LCM / LCM-LoRA + TensorRT}.
  \item Current default safe boundary: \textbf{per-request}.
  \item Strongest imported candidate boundary: \textbf{denoising step}.
  \item Main profiling question: can the runtime make denoising cheap enough, memory-stable enough, and state-aware enough for edge execution under bounded GPU availability?
  \item Current deployment focus:
  \begin{itemize}
    \item \textbf{TensorRT},
    \item \textbf{Triton},
    \item optional \textbf{Jetson Platform Services / MIG} as deployment or isolation scaffolds.
  \end{itemize}
\end{itemize}
\slidecite{Latent Diffusion Models (2022); LCM (2023); Triton Stable Diffusion tutorial; TensorRT docs.}
\end{frame}

\begin{frame}[t,shrink=9]{Slide 6. Computer Vision}
\small
\begin{itemize}
  \item Category type: one-shot, activation-heavy vision inference.
  \item Representative packet-1 anchor: \textbf{RTMPose-m + TensorRT}.
  \item This category currently covers three completed packet families:
  \begin{itemize}
    \item \textbf{Image segmentation}: workload-aligned anchor \textbf{SegFormer / DeepLabV3+}, deployment anchor \textbf{YOLO segmentation}.
    \item \textbf{Pose estimation}: workload-aligned anchor \textbf{RTMPose with ViTPose reference}, deployment anchor \textbf{RTMPose-m / RTMPose-s}.
    \item \textbf{Super resolution}: workload-aligned anchor \textbf{Real-ESRGAN / SwinIR / HAT}, deployment anchor \textbf{Real-ESRGAN}.
  \end{itemize}
  \item Common runtime pattern: \textbf{TensorRT-first}.
  \item Common profiling question: can the runtime validate any useful \textbf{coarse layer-group / tile boundary}, or does it fall back to \textbf{per-inference}?
\end{itemize}
\slidecite{SegFormer (2021); RTMPose (2023); Real-ESRGAN (2021); USHER (2024); Orion (2024); PPipe (2025).}
\end{frame}

\begin{frame}[t]{Slide 7. Voice-related}
\small
\begin{itemize}
  \item This category is \textbf{split first} before profiling:
  \begin{itemize}
    \item streaming ASR,
    \item offline ASR,
    \item speech translation,
    \item TTS.
  \end{itemize}
  \item Packet-1 branch selected: \textbf{streaming ASR}.
  \item Representative packet-1 anchor: \textbf{Emformer-style streaming ASR + ONNX Runtime Mobile}.
  \item Current default safe boundary: \textbf{endpoint / segment}.
  \item Strongest imported candidate boundary: \textbf{audio chunk}.
  \item Main profiling question: can the runtime carry streaming state efficiently enough on embedded hardware without overstating chunk-level reclaim/resume semantics?
\end{itemize}
\slidecite{Emformer (2020); sherpa-onnx streaming docs; ONNX Runtime Mobile docs.}
\end{frame}

\begin{frame}[t]{Slide 8. Translation}
\small
\begin{itemize}
  \item Category type: \textbf{encoder-decoder seq2seq} with extra fixed \textbf{encoder-output state} in addition to decoder KV.
  \item Representative packet-1 anchor: \textbf{M2M100-418M + CTranslate2}.
  \item Workload-aligned scientific reference: \textbf{NLLB-200 distilled-600M}.
  \item Current default safe boundary: \textbf{per-request}.
  \item Strongest imported candidate boundary: \textbf{decoder-step}.
  \item Main profiling question: can the runtime manage both iterative decoding and the extra persistent encoder-side tensors under the same lease and memory budget?
\end{itemize}
\slidecite{M2M-100 (2020); NLLB (2022); CTranslate2 docs; vLLM / Orca / DistServe as imported mechanisms.}
\end{frame}

\begin{frame}[t]{Slide 9. Volume Rendering}
\small
\begin{itemize}
  \item Category type: scene-state-aware rendering around \textbf{ray-batch} work units.
  \item Representative packet-1 anchor: \textbf{Instant-NGP + tiny-cuda-nn}.
  \item Current best candidate boundary: \textbf{ray-batch}.
  \item Main profiling question: can the system expose useful execution units and compact scene state, even though most primaries are accelerators rather than full serving runtimes?
  \item Current deployment focus:
  \begin{itemize}
    \item \textbf{native CUDA / tiny-cuda-nn},
    \item candidate wrappers like \textbf{Triton + K3s/Kubernetes},
    \item challengers such as \textbf{TensoRF} and \textbf{DirectVoxGO}.
  \end{itemize}
\end{itemize}
\slidecite{Instant-NGP (2022); NerfAcc (2023); TensoRF (2022); DirectVoxGO (2022).}
\end{frame}

\begin{frame}[t]{Slide 10. Current Selection Goal}
\small
\begin{itemize}
  \item For each \textbf{category}, I am identifying the \textbf{most suitable runtime or runtime-system stack}, not just the most interesting model.
  \item The runtime will be preferred if it already satisfies the profiling scorecard, or if it can be modified with bounded effort to satisfy it.
  \item The output of this work is a \textbf{category-to-runtime mapping} that supports later prototype design and integration.
  \item In short: the aim is to select, per category, the stack that is most realistic for safe and efficient inference on shared RAN-edge infrastructure.
\end{itemize}
\slidecite{profile\_doc.md and completed packet set under progress/.}
\end{frame}

\begin{frame}[allowframebreaks]{Selected References}
\tiny
\begin{enumerate}
  \item Michele Polese, Niloofar Mohamadi, Salvatore D'Oro, Leonardo Bonati, and Tommaso Melodia. \emph{Beyond Connectivity: An Open Architecture for AI-RAN Convergence in 6G}. arXiv preprint arXiv:2507.06911, 2025.
  \item Chenyuan Feng et al. \emph{AI-RAN: The Pathway to Future Wireless Networks}. Journal of Information and Intelligence, 2026.
  \item Sunghyun Jin et al. \emph{End-to-End Coordination of RAN and Edge Server for Latency-Critical Inference Serving over Cellular Networks}. ACM CoNEXT, 2025.
  \item Woosuk Kwon et al. \emph{Efficient Memory Management for Large Language Model Serving with PagedAttention}. SOSP, 2023.
  \item Gyeong-In Yu et al. \emph{Orca: A Distributed Serving System for Transformer-Based Generative Models}. OSDI, 2022.
  \item Yinmin Zhong et al. \emph{DistServe: Disaggregating Prefill and Decoding for Goodput-optimized Large Language Model Serving}. OSDI, 2024.
  \item Gu et al. \emph{SpotServe: Serving Generative Large Language Models on Preemptible Instances}. ASPLOS, 2024.
  \item Rombach et al. \emph{High-Resolution Image Synthesis With Latent Diffusion Models}. CVPR, 2022.
  \item Luo et al. \emph{Latent Consistency Models: Synthesizing High-Resolution Images with Few-Step Inference}. 2023.
  \item Xue et al. \emph{RTMPose: Real-Time Multi-Person Pose Estimation}. 2023.
  \item Liang et al. \emph{SwinIR: Image Restoration Using Swin Transformer}. ICCVW, 2021.
  \item Wang et al. \emph{Real-ESRGAN: Training Real-World Blind Super-Resolution with Pure Synthetic Data}. ICCVW, 2021.
  \item Fan et al. \emph{Emformer: Efficient Memory Transformer Based Acoustic Model for Low Latency Streaming Speech Recognition}. 2020.
  \item Fan et al. \emph{No Language Left Behind (NLLB)}. 2022.
  \item Fan et al. \emph{Beyond English-Centric Multilingual Machine Translation}. 2022.
  \item Thomas M{"u}ller et al. \emph{Instant Neural Graphics Primitives with a Multiresolution Hash Encoding}. SIGGRAPH, 2022.
\end{enumerate}
\end{frame}

\end{document}
